% ==========================================
% Kapitel 12: Holomorphe Funktionen in Kreisringen und Laurentreihen
% ==========================================
\setcounter{section}{11} % Setzt Zähler, falls du Kapitel 11 bereits hast
\section{Holomorphe Funktionen in Kreisringen und Laurentreihen}



\begin{defi}{12.1 Kreisring (Annulus)}
Seien $0 \le r < R \le \infty$ und $z_0 \in \C$. Der \textbf{Kreisring} (Annulus) um $z_0$ mit innerem Radius $r$ und äußerem Radius $R$ ist definiert als:
\[
    A_{r,R}(z_0) := \{z \in \C : r < |z-z_0| < R\}
\]
Spezialfälle:
\begin{itemize}
    \item $A_{0,R}(z_0) = B_R(z_0) \setminus \{z_0\}$ (punktierte Kreisscheibe)
    \item $A_{0,\infty} = \C \setminus \{z_0\}$
    \item $A_{r,R}(z_0) = B_R(z_0) \setminus \overline{B_r(z_0)}$
\end{itemize}
\end{defi}

\begin{defi}{12.2 Laurentreihe}
Eine \textbf{Laurentreihe} um $z_0$ ist eine Reihe der Form
\[
    \sum_{n=-\infty}^\infty a_n (z-z_0)^n
\]
Sie besteht aus:
\begin{itemize}
    \item \textbf{Hauptteil}: $\sum_{n=-\infty}^{-1} a_n (z-z_0)^n$
    \item \textbf{Nebenteil (Regulärer Teil)}: $\sum_{n=0}^\infty a_n (z-z_0)^n$
\end{itemize}
Die Reihe konvergiert in $z_1$, wenn sowohl Hauptteil als auch Nebenteil konvergieren.
\end{defi}

\begin{satz}{12.3 Holomorphie von Laurentreihen}
Seien $a_n \in \C$ und die Reihe $\sum_{n=-\infty}^\infty a_n (z-z_0)^n$ konvergiere in $z_1$ und $z_2$ mit $|z_1-z_0| < |z_2-z_0|$. Dann konvergiert die Laurentreihe absolut und lokal gleichmäßig auf dem Kreisring $A_{|z_1-z_0|, |z_2-z_0|}(z_0)$ und stellt eine holomorphe Funktion $f$ dar. Die Laurentkoeffizienten sind eindeutig bestimmt durch:
\[
    a_n = \frac{1}{2\pi\I} \int_{\partial B_\rho(z_0)} \frac{f(w)}{(w-z_0)^{n+1}} \, dw
\]
für jedes $\rho \in (r, R)$.
\end{satz}

\begin{satz}{12.4 Cauchysche Ungleichungen für Laurentreihen}
Sei $f(z) = \sum_{n=-\infty}^\infty a_n (z-z_0)^n$ in $A_{r,R}(z_0)$ konvergent. Dann gilt für jedes $\delta \in (r, R)$:
\[
    |a_n| \le \delta^{-n} \sup_{z \in \partial B_\delta(z_0)} |f(z)| \quad \forall n \in \Z
\]
\end{satz}

\begin{satz}{12.7 Pfadunabhängigkeit im Kreisring}
Sei $f: A_{r,R}(z_0) \to \C$ holomorph. Dann ist das Integral
\[
    \int_{\partial B_\delta(z_0)} f(z) \, \dz
\]
unabhängig von $\delta \in (r, R)$.
\end{satz}